\section{Extended Archival Ledgers of Inquilinaje, Agrarian Labor Petitions, and Union Density (1932--1975)}
\label{app:historical_labor_ledgers}

This appendix provides historical ledgers supporting Section~\ref{sec:historical_background} regarding rural labor mobilization, the 1947--48 restrictive legislation, urban land seizures, and the strike wave of the Long Sixties.
% SEC51-VERIFY-14: Audit inherited episode annotations and series definitions
% against source passages before submission; counts are unchanged in this pass.

\subsection{Annual Archival Ledger of Agricultural Labor Petitions (1932--1950)}
\label{app:petitions_ledger}

Table~\ref{tab:app_petitions_full} presents agricultural labor petitions (\textit{Pliegos de Peticiones}) registered by the Labor Department (\textit{Dirección General del Trabajo}), including the postwar crest and subsequent contraction. The restrictive legislation comprised two distinct statutes: agricultural-union Law 8,811 of 1947 and Law 8,987 of 1948, the \textit{Ley Maldita} \citep{Chile1947,Chile1948}. The petition series draws on the compilations of \citet{MorrisOyaneder1962} and the archival research of \citet{Loveman1971,Loveman1976}.

\begin{table}[htbp]
\centering
\footnotesize
\caption{Annual Agricultural Labor Petitions Registered by the Labor Department (Chile, 1932--1950)}
\label{tab:app_petitions_full}
\begin{tabular}{lcp{9.5cm}}
\toprule
\textbf{Year} & \textbf{Petitions Registered} & \textbf{Institutional Context and State Response} \\
\midrule
1932 & 5 & Re-constitutionalization under Arturo Alessandri Palma; initial rural discontent \\
1933 & 8 & Enactment of 1931 Labor Code enforcement decrees; agricultural unionization barred \\
1934 & 12 & Ranquil peasant uprising brutally suppressed by Carabineros \\
1935 & 10 & Landlord resistance to minimum wage decrees \\
1936 & 14 & Creation of the Confederation of Workers of Chile (CTCH) \\
1937 & 18 & Parliamentary elections; rural agitation \\
1938 & 26 & Popular Front electoral victory of Pedro Aguirre Cerda \\
1939 & 171 & Mass unionization push; peasant mobilization breakthrough \\
1940 & 199 & Ministerial suspension circular freezes processing of rural petitions \\
1941 & 82 & Aguirre Cerda death; temporary containment of agrarian conflict \\
1942 & 64 & Election of Juan Antonio Ríos; wartime price controls \\
1943 & 58 & CTCH National Congress; rank-and-file rural pressure resumes \\
1944 & 71 & Radical Party agrarian debates \\
1945 & 95 & End of World War II; urban cost-of-living strikes diffuse into countryside \\
1946 & 284 & Election of Gabriel González Videla with Communist Party (PCCh) participation \\
\textbf{1947} & \textbf{383} & \textbf{All-time historical peak; 31.1\% of all national labor disputes in Chile} \\
\textbf{1948} & \textbf{52} & \textbf{Law 8,987 (\textit{Ley Maldita}); PCCh banned. Agricultural-union Law 8,811 dates from 1947.} \\
1949 & 19 & Mass purges of rural union leadership; Pisagua internment camp opens \\
\textbf{1950} & \textbf{11} & \textbf{Near-total collapse of agricultural collective petitions (-97.1\% from 1947)} \\
\bottomrule
\end{tabular}
\begin{minipage}{\textwidth}
\vspace{0.3em}
\footnotesize
\textit{Source:} Primary archival records from the \textit{Dirección General del Trabajo}, compiled in \citet{MorrisOyaneder1962} and \citet{Loveman1971,Loveman1976}.
\end{minipage}
\end{table}

\subsection{Strike Typology and Legal Status (1960--1975)}
\label{app:strike_typology}

Table~\ref{tab:app_strikes_cliolab} reproduces the complete strike series from Clio-Lab PUC (Sheet 7.7), highlighting the explosion of illegal strikes during the Long Sixties and the Unidad Popular.

\begin{table}[htbp]
\centering
\footnotesize
\setlength{\tabcolsep}{4.5pt}
\caption{Strike Typology, Illegality Rates, and Agricultural Union Membership (1960--1975)}
\label{tab:app_strikes_cliolab}
\begin{tabular}{lcccccc}
\toprule
\textbf{Year} & \shortstack{\textbf{Total}\\\textbf{Strikes}} & \shortstack{\textbf{Legal}\\\textbf{Strikes}} & \shortstack{\textbf{Illegal}\\\textbf{Strikes}} & \shortstack{\textbf{\%}\\\textbf{Illegal}} & \shortstack{\textbf{Total}\\\textbf{Strikers}} & \shortstack{\textbf{Peasant}\\\textbf{Unionists}} \\
\midrule
1960 & 259 & 88 & 171 & 66.0\% & 89,314 & -- \\
1961 & 262 & 95 & 167 & 63.7\% & 112,450 & -- \\
1962 & 401 & 112 & 289 & 72.1\% & 124,100 & -- \\
1963 & 413 & 101 & 312 & 75.5\% & 118,900 & -- \\
1964 & 676 & 142 & 534 & 79.0\% & 138,500 & 1,652 \\
1965 & 723 & 131 & 592 & 81.9\% & 182,354 & 2,118 \\
1966 & 1,073 & 186 & 887 & 82.7\% & 195,400 & 10,647 \\
1967 & 1,114 & 168 & 946 & 84.9\% & 225,120 & 47,483 \\
1968 & 1,138 & 172 & 966 & 84.9\% & 276,450 & 76,356 \\
1969 & 1,277 & 145 & 1,132 & 88.6\% & 356,800 & 104,334 \\
1970 & 1,853 & 179 & 1,674 & 90.3\% & 416,500 & 143,142 \\
1971 & 2,709 & 181 & 2528 & 93.3\% & 485,300 & 232,160 \\
1972 & 3,037 & 118 & 2,919 & 96.1\% & 542,800 & 242,500 \\
1973 & 1,580 & 62 & 1,518 & 96.1\% & 312,400 & 250,740 \\
1974 & 0 & 0 & 0 & -- & 0 & -- \\
1975 & 0 & 0 & 0 & -- & 0 & -- \\
\bottomrule
\end{tabular}
\begin{minipage}{\textwidth}
\vspace{0.3em}
\footnotesize
\textit{Source:} Clio-Lab PUC (Sheet 7.7, Personas). Zero strikes recorded in 1974--1975 reflect the total ban on industrial disputes, strikes, and collective bargaining under military Decree Law 198.
\end{minipage}
\end{table}

\subsection{Trade Union Affiliation and Union Density Ledger (1932--1975)}
\label{app:union_density_ledger}

Table~\ref{tab:union_density_osorio_polanco} reports the complete archival series on trade union affiliation, total employed workforce, average union size, and union density over employed workers reconstructed by \citet{OsorioPolanco2021}, documenting the 1948--1958 stagnation under the \textit{Ley Maldita} and the historic mobilization peak under the Unidad Popular.

\begin{table}[htbp]
\centering
\footnotesize
\setlength{\tabcolsep}{3.5pt}
\renewcommand{\arraystretch}{0.90}
\caption{Trade Union Organization, Density over Employed Workers, and Average Union Size (1932--1975)}
\label{tab:union_density_osorio_polanco}
\begin{tabularx}{\textwidth}{lcccccX}
\toprule
\textbf{Year} & \textbf{Unions} & \textbf{Affiliated} & \textbf{Employed} & \textbf{Avg Size} & \textbf{Density (\%)} & \textbf{Accumulation Regime \& Institutional Era} \\
\midrule
1932 & 421 & 54,801 & 1,165,000 & 130.2 & 4.70\% & Post-Depression trough; Alessandri restoration \\
1935 & 669 & 83,262 & 1,553,000 & 124.5 & 5.36\% & Early ISI expansion under 1931 Labor Code \\
1938 & 932 & 125,972 & 1,676,000 & 135.2 & 7.52\% & Popular Front victory (Aguirre Cerda) \\
1939 & 1,687 & 173,438 & 1,709,000 & 102.8 & 10.15\% & Democratic surge; enterprise union proliferation \\
1940 & 1,888 & 162,297 & 1,746,000 & 86.0 & 9.30\% & Secret circular freezes peasant unionization \\
1941 & 1,977 & 208,779 & 1,763,000 & 105.6 & 11.84\% & Industrialization acceleration under CORFO \\
1944 & 1,652 & 246,221 & 1,852,000 & 149.0 & 13.29\% & Wartime union consolidation; CTCH congress \\
1947 & 1,831 & 263,085 & 1,926,000 & 143.7 & \textbf{13.66\%} & Pre-\textit{Ley Maldita} peak of labor mobilization \\
1948 & 1,857 & 263,676 & 2,029,000 & 142.0 & \textbf{13.00\%} & Enactment of \textit{Ley Maldita} (Laws 8,987 \& 8,811) \\
1950 & 1,907 & 260,143 & 2,087,000 & 136.4 & 12.46\% & Repression of left leadership; union growth freezes \\
1952 & 1,997 & 284,418 & 2,202,000 & 142.4 & 12.92\% & Ibáñez del Campo populist election \\
1955 & 2,177 & 305,192 & 2,209,000 & 140.2 & 13.82\% & Klein-Saks mission wage freeze \\
1958 & 1,894 & 276,346 & 2,297,000 & 145.9 & 12.03\% & Repeal of \textit{Ley Maldita}; density stagnation \\
1960 & 1,770 & 272,966 & 2,313,000 & 154.2 & 11.80\% & Post-repeal realignment; Jorge Alessandri \\
1964 & 1,863 & 278,980 & 2,497,000 & 149.7 & \textbf{11.17\%} & Frei Montalva election; rural baseline \\
1966 & 2,882 & 369,507 & 2,604,000 & 128.2 & 14.19\% & Peasant organizing begins (\textit{Promoción Popular}) \\
1967 & 3,426 & 442,650 & 2,681,000 & 129.2 & 16.51\% & Laws 16,625 (Peasant Unions) \& 16,640 (Agrarian Reform) \\
1968 & 3,889 & 500,447 & 2,714,000 & 128.7 & 18.44\% & Agrarian radicalization \& urban strike waves \\
1969 & 4,228 & 541,967 & 2,734,000 & 128.2 & 19.82\% & Pre-electoral polarization; \textit{El Tacnazo} \\
1970 & 4,581 & 627,664 & 2,743,000 & 137.0 & \textbf{22.88\%} & Election of Salvador Allende (Unidad Popular) \\
1971 & 5,212 & 782,494 & 2,829,000 & 150.1 & 27.66\% & Reactivation boom; massive unionization surge \\
1972 & 6,118 & 855,404 & 2,932,000 & 139.8 & \textbf{29.17\%} & Distributive peak; Bosses' Lockout (\textit{Paro de Octubre}) \\
1973 & 6,692 & 939,319 & 2,976,000 & 140.4 & \textbf{31.56\%} & All-time historical zenith of union density \\
1974 & 7,069 & 947,093 & 2,907,000 & 134.0 & 32.58\% & Military Junta halts collective bargaining \\
1975 & 7,181 & 940,810 & 2,683,000 & 131.0 & 35.07\% & Shock therapy depression; employment contraction \\
\bottomrule
\end{tabularx}
\begin{minipage}{\textwidth}
\vspace{0.3em}
\footnotesize
\textit{Source:} \citet{OsorioPolanco2021}, \textit{Repositorio de Estadísticas Sindicales (RES)}, sheet \texttt{Empalme}. Average union size is $\text{Afiliación} / \text{Sindicatos}$. Union density is $\text{Afiliación} / \text{Ocupados}$.
\end{minipage}
\end{table}


\subsection{Urban Land Seizures in Santiago and Nationwide (1967--1971)}
\label{app:land_seizures_ledger}

Table~\ref{tab:urban_land_seizures_murphy} details the exponential expansion of squatter land seizures (\textit{tomas de terreno}) in Greater Santiago and nationwide from 1967 through 1971, documenting the spatio-temporal displacement of peripheral class conflict from the countryside into urban social reproduction \citep{Murphy2015}.

\begin{table}[htbp]
\centering
\small
\caption{Urban Land Seizures in Santiago and Nationwide (1967--1971)}
\label{tab:urban_land_seizures_murphy}
\begin{tabularx}{\textwidth}{l c c c X}
\toprule
\textbf{Year} & \shortstack{\textbf{Santiago Land}\\\textbf{Seizures}} & \shortstack{\textbf{Chile Total}\\\textbf{Land Seizures}} & \shortstack{\textbf{Santiago}\\\textbf{Share (\%)}} & \textbf{Historical Phase} \\
\midrule
1967 & 13 & 8 & --- & Agrarian Reform \& Unionization Laws promulgated \\
1968 & 4 & 35 & 11.4\% & Consolidation of \textit{pobladores} committees \\
1969 & 23 & 220 & 10.5\% & Radicalization; Puerto Montt massacre \\
1970 & 103 & 560 & \textbf{18.4\%} & \textbf{Election of Salvador Allende (UP)} \\
1971 & 350 & --- & --- & \textbf{Reactivation boom \& urban mobilization} \\
\bottomrule
\end{tabularx}
\begin{minipage}{\textwidth}
\vspace{0.3em}
\footnotesize
\textit{Source:} \citet{Murphy2015}, p.~58. \textit{Note:} Demonstrates the exponential growth of the \textit{pobladores} movement and the spatio-temporal displacement of the class struggle into urban informality and collective social reproduction. In 1967, Santiago reported 13 municipal seizures prior to full national reporting standardization.
\end{minipage}
\end{table}


\subsection{Productive Capital Stock, Profit Rate, and Net Accumulation Framework (1960--1975)}
\label{app:profit_accumulation_ledger}

This subsection formalizes the balance-sheet accounting, Generalized Perpetual Inventory Method (GPIM) depreciation formulas, and Weisskopf-Marxian decomposition equations supporting Section~4 §4.10 regarding the Marxian profit squeeze and net capital accumulation.

\paragraph{Productive Fixed Capital Stock Harmonization.}
Following the capital stock harmonization methodology established in \citet{Polanco2026}, the capital stock denominator is restricted strictly to productive fixed assets, encompassing Machinery and Equipment ($ME$) and Non-Residential Construction ($NRC$):
\begin{equation}
	K_t \equiv K_{\text{ME}, t} + K_{\text{NRC}, t}
	\label{eq:k_definition_app}
\end{equation}
Residential Construction ($RC$) is explicitly excluded because housing represents social overhead infrastructure and durable consumption yielding utility rather than direct surplus-value producing capital.

\paragraph{The Weisskopf Decomposition.}
Drawing on the Classical-Marxian and Post-Keynesian traditions \citep{Weisskopf1979,Foley1982,Basu2022}, the aggregate profit rate on productive capital ($r_t$) at current replacement cost decomposes into the profit share, capacity utilization, and potential capital productivity:
\begin{equation}
	r_t \equiv \frac{\Pi_t}{P_{K, t} K_t} = (1 - \omega_t) \cdot \frac{Y_t}{P_{K, t} K_t} = (1 - \omega_t) \cdot \mu_t \cdot \sigma_t
	\label{eq:profit_rate_weisskopf_app}
\end{equation}
where $\Pi_t \equiv (1 - \omega_t) Y^{\text{nom}}_t$ is the nominal mass of aggregate profits (surplus value), $\omega_t$ is Astorga's \citeyearpar{Astorga2023} aggregate wage share, $1 - \omega_t$ is the profit share, $\mu_t \equiv Y_t / Y^p_t$ is structural capacity utilization (estimated via Cointegrating Multivariate Polynomial Regression IM-OLS in \citealt{Polanco2026}), $\sigma_t \equiv Y^p_t / (P_{K, t} K_t)$ is potential capital productivity, and $Y_t / (P_{K, t} K_t) = \mu_t \sigma_t$ is the actual output-capital ratio at current replacement cost. In standard macroeconomic notation, I denote the aggregate profit rate by $r$, reserving $u$ for the open unemployment rate, $\mu$ for capacity utilization, and $i$ for nominal interest rates.

\paragraph{Implicit Physical Depreciation and Net Accumulation.}
Under perpetual inventory balance-sheet accounting, the net expansion of the capital stock requires that the rate of capital accumulation be discounted from the implicit annual depreciation of net stocks. Implicit annual physical depreciation ($D_t$) and the implicit physical depreciation rate ($\delta_t$) are defined as:
\begin{equation}
	D_t \equiv K_{t-1} + I_t - K_t \qquad \text{and} \qquad \delta_t \equiv \frac{D_t}{K_{t-1}}
	\label{eq:implicit_depreciation_app}
\end{equation}
where $I_t \equiv I_{\text{ME}, t} + I_{\text{NRC}, t}$ is real gross fixed investment in productive assets. 

The Net Rate of Capital Accumulation ($g^n_{K, t}$) and the Re-capitalization Propensity out of surplus value ($\chi_t \equiv I^{\text{net}}_t / \Pi_t$) are governed by:
\begin{equation}
	g^n_{K, t} \equiv \frac{I_t - D_t}{K_{t-1}} = g^{\text{gross}}_{K, t} - \delta_t = \chi_t \cdot r_t = \chi_t \cdot (1 - \omega_t) \cdot \mu_t \cdot \sigma_t
	\label{eq:net_accumulation_rate_app}
\end{equation}
Equation~\eqref{eq:net_accumulation_rate_app} formalizes the Cambridge-Marxian accumulation identity: the net growth rate of productive capital equals the product of the re-capitalization propensity ($\chi$) and the aggregate profit rate ($r$).

Table~\ref{tab:profit_rate_accumulation_1960_1975} reports the complete numerical trajectory of these series from 1960 to 1975.

\begin{table}[htbp]
\centering
\footnotesize
\setlength{\tabcolsep}{3.5pt}
\caption{The Distributive Trajectory, Capacity Utilization, Profit Rate, and Net Capital Accumulation (1960--1975)}
\label{tab:profit_rate_accumulation_1960_1975}
\resizebox{\textwidth}{!}{%
\begin{tabular}{lcccccccl}
\toprule
\textbf{Year} & $\bm{\omega}$ \textbf{(\%)} & $\bm{u}$ \textbf{(\%)} & $\bm{\mu}$ \textbf{(\%)} & $\bm{Y/K}$ & $\bm{r}$ \textbf{(\%)} & $\bm{\chi}$ & $\bm{g^n_K}$ \textbf{(\%)} & \textbf{Regime Dynamic} \\
\midrule
1960 & 60.9\% & 6.6\% & 86.9\% & 0.2580 & 10.08\% & 0.421 & 4.44\% & Pre-reform stabilization baseline \\
1961 & 65.9\% & 7.5\% & 87.3\% & 0.2424 & 8.27\% & 0.534 & 4.62\% & Early ISI expansion \\
1962 & 59.4\% & 7.4\% & 87.9\% & 0.2593 & 10.52\% & 0.421 & 4.63\% & Exchange rate devaluation \\
1963 & 55.8\% & 7.0\% & 89.8\% & 0.2670 & 11.81\% & 0.423 & 5.25\% & Pre-electoral demand expansion \\
1964 & 67.7\% & 6.5\% & 88.4\% & 0.2525 & 8.16\% & 0.563 & 4.82\% & Frei Montalva election; wage surge \\
1965 & 63.8\% & 5.9\% & 86.7\% & 0.2374 & 8.59\% & 0.420 & 3.75\% & Initial anti-inflationary program \\
1966 & 55.3\% & 5.6\% & 93.4\% & 0.2450 & 10.94\% & 0.323 & 3.66\% & Output peak under Frei \\
1967 & 57.9\% & 4.2\% & 93.5\% & 0.2410 & 10.14\% & 0.356 & 3.75\% & Agrarian reform promulgated; profit squeeze \\
1968 & 53.9\% & 4.4\% & 93.7\% & 0.2452 & 11.29\% & 0.337 & 3.96\% & Drought and fiscal austerity \\
1969 & 50.1\% & 5.0\% & 94.1\% & 0.2560 & 12.79\% & 0.288 & 3.82\% & Pre-electoral stagflation; \textit{El Tacnazo} \\
1970 & \textbf{50.9\%} & 5.9\% & 92.8\% & 0.2491 & \textbf{12.23\%} & 0.332 & \textbf{4.23\%} & Allende election; inherited slack baseline \\
1971 & \textbf{56.2\%} & 5.2\% & \textbf{98.4\%} & 0.2427 & \textbf{10.62\%} & 0.304 & \textbf{3.34\%} & Idle capacity mobilized; wage expansion \\
1972 & \textbf{69.0\%} & \textbf{4.0\%} & \textbf{96.0\%} & 0.2121 & \textbf{6.57\%} & \textbf{0.277} & \textbf{1.85\%} & Profit squeeze peak; Bosses' Lockout \\
1973 & \textbf{54.5\%} & 4.8\% & 89.5\% & 0.2004 & \textbf{9.12\%} & \textbf{0.184} & \textbf{1.71\%} & Coup d'état; DL 522 price deregulation \\
1974 & \textbf{39.4\%} & 9.1\% & 88.8\% & 0.1762 & \textbf{10.67\%} & 0.218 & \textbf{2.38\%} & Junta wage compression; union ban \\
1975 & 49.3\% & \textbf{18.0\%} & \textbf{76.6\%} & 0.1289 & 6.54\% & 0.183 & \textbf{1.21\%} & Chicago shock therapy depression \\
\bottomrule
\end{tabular}%
}
\begin{minipage}{0.96\textwidth}
\vspace{0.3em}
\footnotesize
\textit{Source:} Author's calculations based on \citet{Astorga2023}, Clio-Lab PUC / \citet{DiazLudersWagner2016}, and \citet{Polanco2026}. Following \citet{Polanco2026}, productive net capital stock ($K_t \equiv K_{\text{ME}, t} + K_{\text{NRC}, t}$) strictly excludes Residential Construction ($RC$). Structural capacity utilization ($\mu \equiv Y/Y^p$) is estimated via the CMPR IM-OLS cointegrating framework of \citet{Polanco2026} (\texttt{Capacity-Utilization-US\_Chile}); $u$ is the open unemployment rate. Implicit annual depreciation is $D_t \equiv K_{t-1} + I_t - K_t$ with depreciation rate $\delta_t \equiv D_t / K_{t-1}$. The aggregate profit rate satisfies the Weisskopf-Marxian decomposition $r_t \equiv (1 - \omega_t) Y_t / (P_{K, t} K_t) = (1 - \omega_t) \mu_t \sigma_t$, where $\sigma_t \equiv Y^p_t / K_t$ is potential capital productivity. Net capital accumulation satisfies the classical Cambridge-Marxian identity $g^n_{K, t} \equiv (I_t - D_t) / K_{t-1} = \chi_t \cdot r_t$, where $\chi_t \equiv I^{\text{net}}_t / \Pi_t$ is the capitalist re-capitalization propensity out of surplus value $\Pi_t \equiv (1 - \omega_t) Y_t$.
\end{minipage}
\end{table}

